% problem-000014 3 磷氧化物结构比较
\section*{3. 磷氧化物结构比较}
\textit{下列为分问。}

\subsection*{3-1}
好奇⼼是科学发展的内在动⼒之⼀。 $\mathrm { P _ { 2 } O _ { 3 } }$ 和 $\mathrm { P _ { 2 } O _ { 5 } }$ 是两种经典的化合物，其分⼦结构已经确定。⾃然⽽然会有如下问题：是否存在磷氧原⼦⽐介于⼆者之间的化合物？由此出发，化学家合成并证实了这些中间化合物的存在。

\subsection*{3-1}
-1 写出这些中间化合物的分⼦式。

\subsection*{3-1}
-2 画出其中具有2重旋转轴的分⼦的结构图。根据键⻓不同，将P—O键分组并⽤阿拉伯数字标出（键⻓相同的⽤同⼀个数字标识）。⽐较键⻆ $\angle 0 { \mathrm { - } } \mathrm { P } ( \mathrm { V } ) { \mathrm { - } } 0$ 和 $\angle 0 { \mathrm { - P ( I I I ) { - } O } }$ 的⼤⼩。

\subsection*{3-2}
$\mathrm { N H _ { 3 } }$ 分⼦独⽴存在时H—N—H键⻆为106.7°。右图是 $\mathrm { [ Z n ( N H _ { 3 } ) _ { 6 } ] } ^ { \prime }$ 2+离⼦的部分结构以及 $_ \mathrm { H - N - H }$ 键⻆的测量值。解释配合物中H—N—H键⻆变为109.5°的原因。

$
\mathrm{Zn} ^ {2 +} \xleftarrow {1 0 9 . 5 ^ {\circ}} \begin{array}{c} \mathrm{H} \\ \mathrm{N} \\ \mathrm{H} \end{array} \begin{array}{c} 1 0 9. 5 ^ {\circ} \\ \mathrm{H} \\ \mathrm{H} \end{array}
$

\subsection*{3-3}
量⼦化学计算预测未知化合物是现代化学发展的途径之⼀。2016年2⽉有⼈通过计算预⾔铁也存在四氧化物，其分⼦构型是四⾯体，但该分⼦中铁的氧化态是+6⽽不是 $+ 8 _ { \mathsf { c } }$

\subsection*{3-3}
-1 写出该分⼦中铁的价电⼦组态。

\subsection*{3-3}
-2 画出该分⼦结构的⽰意图（⽤元素符号表⽰原⼦，⽤短线表⽰原⼦间的化学键）。
